\chapter{Marco Teórico}



\section{Modelo del panel solar}

El modelo de un solo diodo equivalente de una célula de un panel solar con $n_S$ diodos en serie se puede 
observar en el siguiente circuito \cite{ModeloPVDatasheet}: 

% Circuito equivalente PV

\begin{figure}[h]
    \centering
    \begin{circuitikz}[american, scale=1.2, every node/.style={scale=1}]

        % Fuente de corriente Iph
        \draw (0,0) to[isource, l_=$I_{ph}$] (0,4);

        % Cable superior desde la fuente hasta el diodo
        \draw (0,4) -- (2,4);

        % Diodo D con corriente iD
        \draw (2,4) to[D, l_=$D$, i=$i_D$] (2,0);

        % Cable inferior desde el diodo hasta la fuente
        \draw (2,0) -- (0,0);

        % Cable superior desde el diodo hasta Rsh
        \draw (2,4) -- (4,4);

        % Resistencia Rsh con corriente iRsh
        \draw (4,4) to[R, l_=$R_{SH}$, i=$i_{Rsh}$] (4,0);

        % Cable inferior desde Rsh
        \draw (4,0) -- (2,0);

        % Resistencia serie Rs
        \draw (4,4) to[R, l=$R_S$, i=$i$] (7,4);

        % Terminal de salida superior (abierto, con flecha de corriente)
        \draw (7,4) -- (8,4);

        % Cable inferior hasta el terminal de salida
        \draw (4,0) -- (8,0);

        % Etiqueta de voltaje v en el terminal de salida
        \draw (8,4) to[open, v=$v$] (8,0);

    \end{circuitikz}
    \caption{Modelo eléctrico equivalente de una celda solar con resistencia serie $R_S$ y resistencia shunt $R_{SH}$.}
    \label{fig:modelo_celda_solar}
\end{figure}

Aplicando el análisis de circuitos se obtiene que la corriente de salida viene definida por tanto por:

\begin{equation}
i = I_{ph}(I_{rr}) - i_D - i_{R_{SH}} =I_{ph} -I_0(e^{\frac{v+i R_S}{n_S V_t}} - 1)- \frac{v+iR_S}{R_SH}
\label{eq:ModeloDiodoUnico}
\end{equation}

Donde: 
\begin{itemize}
    \item $V_t$ es la tensión térmica de la unión p-n, definida por:
    \begin{equation}
    V_t=\frac{AkT_{STC}}{q}
    \label{eq:TensionTermica}
    \end{equation}
    \item $I_{ph}(I)$ es la fotocorriente determinada en condiciones estándar (STC) con una irradiancia $I_{rr}$.
    \item $I_0$ es la corriente de saturación oscura en condiciones etándar.
    \item $R_S$ es la resistencia en serie de la célula.
    \item $R_{SH}$ es la resistencia Shunt de la célula.
    \item $A$ es el factor de idealidad del diodo.
    \item $T_{STC}[K]$ es la temperatura en las condicones estándar.
    \item $k$ es la constante de Boltzmann, $q$ es la carga del electrón y $n_S$ es el número de células en serie.
\end{itemize}

Nótese que $T_{STC}$, $k$, $q$ y $n_S$ son constantes ya determinadas, por lo que
los únicos parámetros que necesarios para realizar un modelo del comportamiento de la célula solar son $I_{ph}$, $I_0$,
$R_S$,  $R_{SH}$ y $A$.